\section{Integrasi ke Proyek: Token Stream dan Driver}
\label{sec:ch07-integrasi}

\subsection{Driver}

Driver membaca berkas, menginisialisasi lexer, lalu mencetak token
sampai \code{TOK\_EOF}. Implementasinya ditunjukkan pada
\lstref{lst:ch07-main-lex}.

\includecode[
  language=C,
  firstline=1,
  lastline=60,
  firstnumber=1,
  caption={Driver \file{main\_lex.c}},
  label={lst:ch07-main-lex}
]{../../code/subsetc/src/main_lex.c}

Driver ini adalah alat uji, bukan bagian dari compiler final. Pada
milestone berikutnya, driver akan diganti oleh parser yang memanggil
\code{lexer\_next()} secara langsung.

\subsection{Format Keluaran}

Setiap baris keluaran memiliki format:
\[
\text{baris} : \text{kolom} \quad \text{jenis token} \quad
\text{lexeme}.
\]
Format ini memudahkan perbandingan dengan berkas harapan
(\code{.expected}). Contoh keluaran untuk \file{basic.sc}:
\begin{lstlisting}[language=bash]
1:1 INT int
1:5 IDENT main
1:9 LPAREN (
1:10 RPAREN )
1:12 LBRACE {
\end{lstlisting}

\subsection{Pengujian}

Direktori \file{tests/lexer} memuat empat kelompok tes:
\begin{itemize}[leftmargin=*]
  \item \file{basic.sc}: program minimal;
  \item \file{operators.sc}: operator dan tanda baca;
  \item \file{literals.sc}: literal integer, float, char, dan string;
  \item \file{errors.sc}: karakter ilegal, float tidak lengkap, dan
        string tidak ditutup.
\end{itemize}

Setiap berkas \code{.sc} memiliki berkas \code{.expected}. Skrip
\file{run-tests.bat} menjalankan lexer, membandingkan keluaran, dan
melaporkan tes yang gagal.

\subsection{M1: Definisi Selesai}

\begin{checklist}
  \item Lexer tangan dapat dikompilasi tanpa peringatan.
  \item Lexer Flex dapat dihasilkan dan dikompilasi.
  \item Empat kelompok tes tersedia dan lulus.
  \item Kesalahan leksikal menghasilkan diagnostik yang informatif.
  \item Laporan singkat berisi desain, hasil tes, dan kendala.
\end{checklist}

\begin{table}[H]
\centering
\caption{Artefak M1 dan statusnya}
\label{tab:ch07-artefak}
\begin{tabularx}{\textwidth}{@{}l l X@{}}
\toprule
\textbf{Artefak} & \textbf{Status} & \textbf{Keterangan} \\
\midrule
\file{token.h} & selesai & enum dan struktur token \\
\file{lexer.h} & selesai & antarmuka lexer \\
\file{lexer.c} & selesai & scanner tangan \\
\file{lexer\_flex.l} & selesai & scanner Flex \\
\file{error.c} & selesai & pelaporan kesalahan \\
\file{main\_lex.c} & selesai & driver uji \\
\file{tests/lexer} & selesai & empat kelompok tes \\
\bottomrule
\end{tabularx}
\end{table}

\subsection{Catatan Verifikasi}

Lingkungan pengembangan buku ini tidak memiliki GCC, Flex, dan Bison.
Kode pada bab ini disusun dan ditinjau secara statis, sedangkan
keluaran yang diharapkan didokumentasikan pada berkas
\code{.expected}. Mahasiswa harus menjalankan build dan tes pada
lingkungan yang memiliki toolchain C. Jika hasilnya berbeda,
periksa perbedaan versi Flex, penanganan baris baru, dan urutan
aturan.
